% Subproblem 5 (EASY): Lift positivity from one generator to all of T_{2g} via Schur's lemma
%
% CONTEXT: Part of the T_{2g} Hessian proof.
% Working directory: /data/projects/3eacb340-027e-4b71-8de4-67574ba54ea2/QBLYHzbdEjWeyg/proof/
%
% NOTATION:
%   - G = full symmetry group of the cube (order 48, isomorphic to the octahedral group Oh).
%   - V = (direct sum_{i=1}^8 T_{p_i} S^2) / SO(3), a 13-dimensional real representation of G.
%   - T_{2g} is the 3-dimensional irreducible G-subrepresentation of V generated by the
%     deformation v = (v_i) where:
%       v_i = +partial_phi/sin(theta)|_{p_i}  for i in A = {1,2,7,8},
%       v_i = -partial_phi/sin(theta)|_{p_i}  for i in B = {3,4,5,6}.
%   - Q(w,w) = Hess_p log J_8 (w,w) is a G-equivariant quadratic form on V.
%   - From Subproblems 1-4: Q(v, v) = -16*pi^2 * I > 0 for the specific generator v above.
%
% STATEMENT TO PROVE:
%   Q is positive definite on all of T_{2g}: Q(w,w) > 0 for all nonzero w in T_{2g}.
%
% PROOF STRATEGY (Schur's lemma argument):
%   The functional J_8 is invariant under global rotations (SO(3) acts by rotating all p_i
%   simultaneously, and J_8 is rotation-invariant). More importantly, J_8 is invariant
%   under the full symmetry group G of the cubic configuration (relabeling the 8 vertices
%   by cube symmetries). Therefore, the Hessian Q is G-equivariant:
%     Q(g.w, g.w) = Q(w, w)  for all g in G, w in V.
%
%   The restriction of Q to T_{2g} defines a G-equivariant self-adjoint endomorphism
%   L: T_{2g} -> T_{2g} via the formula Q(w, w') = <Lw, w'> for any G-invariant inner product
%   on T_{2g}.
%
%   Since T_{2g} is irreducible (3-dimensional irrep of G), Schur's lemma implies that
%   any G-equivariant endomorphism T_{2g} -> T_{2g} is a scalar multiple of the identity:
%     L = lambda * Id  for some lambda in R.
%
%   Since Q(v,v) = -16*pi^2 * I > 0 for the nonzero generator v, we have lambda > 0.
%   Therefore Q(w,w) = lambda * |w|^2 > 0 for all nonzero w in T_{2g}.
%
% ITEMS TO ADDRESS:
%   (a) State and verify the G-equivariance of Q.
%   (b) State Schur's lemma for real irreducible representations (the real version: for a
%       real irrep over R, any self-adjoint G-equivariant endomorphism is a real scalar
%       multiple of the identity, provided the Schur's lemma division algebra is R).
%       Note: T_{2g} is an absolutely irreducible real representation of G, so Schur applies
%       with the scalar being real.
%   (c) Conclude positive definiteness from lambda > 0.
%
% DELIVERABLE: A clean, concise LaTeX proof (about half a page) of the above.
% Output should be a self-contained LaTeX fragment (no \documentclass).
\end{document}
